% Part: proof-theory
% Chapter: natural-deduction
% Section: grafting

\documentclass[../../../include/open-logic-section]{subfiles}

\begin{document}

\olfileid[tr]{pt}{nat}{gra}

\olsection{\usetoken{P}{proof} Aşılama}

\Log{N1c} ve~\Log{N1i}'deki !!^{proof}s, !!{proof}s{}dır: bunlarda
!!{formula}s~$!A$,~$!B$ varsayımlarından elde edilir. $!B$'nin kendisinin
bir !!a{proof}{}ı
bulunduğunu varsayalım. $!B$'nin !!{proof}{}ıyla $!B$'den $!A$'nın !!{proof}{}ını,
$!B$ varsayımını kullanmayan bir $!A$ !!a{proof}{}ı elde edecek biçimde nasıl
birleştiririz?

Seçeneklerden biri, $!A$'nın !!{proof}{}ından $!B$ varsayımını~\Intro{\lif}
uygulayarak kaldırmaktır. Ardından ortaya çıkan $!B \lif !A$ !!{proof}{}ını
$!B$'nin !!{proof}{}ıyla birleştirerek şunu elde edebiliriz:
\[
\AxiomC{$\Discharge{!B}{x}$}
\DeduceC{$!A$}
\DischargeRule{\Intro{\lif}}{x}
\UnaryInfC{$!B \lif !A$}
\AxiomC{}
\DeduceC{$!B$}
\RightLabel{\Elim{\lif}}
\BinaryInfC{$!A$}
\DisplayProof
\]
Bu kolaydır, fakat biraz çirkindir: $\lif$'i hemen ardından ortadan kaldırmak
için neden önce ortaya sokalım? $!B \lif !A$ üzerinden geçen bu tür bir
``dolambaçtan'' kaçınmak mümkün olmalıdır. (Bunlardan her zaman
kaçınılabildiği, aslında normalleştirme teoreminin içeriğidir.)

\Log{N1c} ve~\Log{N1i}'de bu durumda dolambaçtan gerçekten oldukça kolay
kaçınabiliriz: $!B^x$ varsayımlarını $!B$'nin bütün !!{proof}{}ıyla değiştirmek
yeterlidir. !!{proof}s{}ın başkalarının varsayımlarına bu biçimde
``aşılanması'' oldukça yararlıdır; dolayısıyla bunu bir tanım olarak kayda
geçiriyoruz.

\begin{defn}
$\delta_1$, açık $!B^x$ varsayımından $!A$'nın bir \Log{N1c}-!!{proof}{}ı
(ya da \Log{N1i}-!!{proof}{}ı), $\delta_2$ de $!B$'nin bir
\Log{N1c}-!!{proof}{}ı (\Log{N1i}-!!{proof}{}ı) ise
$\Subst{\delta_1}{\delta_2}{!B^x}$, $\delta_1$ içindeki boşaltılmamış her
$!B^x$ geçişinin $\delta_2$ ile değiştirilmesinin sonucudur.
\end{defn}

$\delta_1$ ile $\delta_2$ aynı etiketle etiketlenmiş farklı varsayım
!!{formula}s{}ine sahip değilse ve $\delta_2$'nin hiçbir açık varsayımı
$\delta_1$'in bir özdeğişkenini içermiyorsa ortaya çıkan
$\Subst{\delta_1}{\delta_2}{!B^x}$ doğru bir !!{proof}{}tır; açık varsayımları
da $\delta_1$'in açık varsayımlarıyla $\delta_2$'ninkilerin birlikte
alınmasıdır. !!{proof}s aşılanırken bu koşullar genellikle sağlanır. Birinci
koşul, $\delta_1$ ile $\delta_2$ daha büyük bir !!{proof}{}ın alt ispatları
olduğunda sağlanır. İkinci koşul sağlanmıyorsa $\delta_1$'in özdeğişkenlerini
sağlanacak biçimde yeniden adlandırabiliriz.

Karşılık gelen bir tanım \Log{N2c} ve~\Log{N2i} için de geçerlidir.

\begin{defn}
$\delta_1$, $x:!B, \Gamma_1 \Sequent !A$'nın bir \Log{N2c}-!!{proof}{}ı
(ya da \Log{N2i}-!!{proof}{}ı), $\delta_2$ de $\Gamma_2 \Sequent !B$'nin bir
\Log{N2c}-!!{proof}{}ı (\Log{N2i}-!!{proof}{}ı) ise
$\Subst{\delta_1}{\delta_2}{x:!B}$, $\delta_1$ içindeki her
$x:!B \Sequent !B$ aksiyomunun $\delta_2$ ile değiştirilmesinin ve
$\Gamma_2$'nin bu aksiyomların altındaki her sekantın bağlamına eklenmesinin
sonucudur.
\end{defn}

\begin{prop}
$\delta_1 \Proves x:!B, \Gamma_1 \Sequent !A$ ve
$\delta_2 \Proves \Gamma_2 \Sequent !B$ ise ve $\delta_1$'in hiçbir
özdeğişkeni $\Gamma_2$ içinde geçmiyorsa
$\Subst{\delta_1}{\delta_2}{x:!B} \Proves \Gamma_1, \Gamma_2 \Sequent !A$.
\end{prop}


\begin{proof}
  $\pheight{\delta_1}$ üzerinde tümevarımla.
\end{proof}

\end{document}
